% =====================================================================
\chapter{Les workflows métier}
\label{ch:workflows}

Chaque workflow est présenté selon la chaîne demandée, en n'employant que les
composants réellement présents dans le code. Quand une étape n'existe pas (pas
de LLM, pas de MCP), elle est simplement absente plutôt que mentionnée pour la
forme.

% =====================================================================
\section{Workflow 1 --- Authentification}

\begin{longtable}{@{}p{2.9cm}p{11.9cm}@{}}
\toprule
\textbf{Étape} & \textbf{Ce qui se passe}\\
\midrule
\endhead
Utilisateur & Saisit e-mail, mot de passe, coche ou non « se souvenir »\\
Frontend & \cls{LoginPage} \arrowto{} \fn{signIn(email, password, remember)} de \cls{AuthContext}\\
API & \route{POST /api/v1/auth/login} --- \textbf{aucune dépendance d'authentification}\\
Backend & \fn{auth\_service.login} \arrowto{} \fn{authenticate\_user} : \texttt{SELECT} sur \tbl{users}, \fn{verify\_password} (bcrypt), contrôle \texttt{is\_active}\\
Résultat & \fn{create\_access\_token(subject=email, role=role.name)} \arrowto{} JWT HS256 avec \texttt{sub}, \texttt{role}, \texttt{exp}\\
Frontend & \fn{tokenStore.set(token, remember)} : \texttt{localStorage} si « se souvenir », sinon \texttt{sessionStorage}. Puis \route{GET /auth/me} pour charger le profil\\
Ensuite & \cls{RequireAuth} laisse passer ; \cls{AppLayout} s'affiche\\
\bottomrule
\end{longtable}

\subsection*{Reprise de session au rechargement}
\cls{AuthProvider} n'accorde pas sa confiance au seul jeton stocké : il appelle
\route{GET /auth/me}. Commentaire du code : « un jeton en stockage ne suffit
pas, il faut que le backend le reconnaisse encore ». En cas d'échec,
\fn{tokenStore.clear()}.

\subsection*{Expiration en cours d'usage}
Toute réponse 401 déclenche \fn{onUnauthorized()} dans \file{lib/api.ts}, câblé
sur \fn{signOut}. Le hook \fn{useResource} \emph{n'affiche pas} l'erreur dans ce
cas : « Un 401 est déjà traité par le handler global : ne pas l'afficher. »

% =====================================================================
\section{Workflow 2 --- Dashboard}

\begin{longtable}{@{}p{2.9cm}p{11.9cm}@{}}
\toprule
\textbf{Étape} & \textbf{Ce qui se passe}\\
\midrule
\endhead
Utilisateur & Arrive sur \route{/storage/dashboard}\\
Frontend & Quatre appels : \route{/overview}, \route{/performance?range=24h}, \route{/activity?limit=6}, \route{/insights}\\
API & \file{api/v1/dashboard.py}, protégé par \fn{require\_role("storage\_engineer","administrator")}. Routes en \texttt{def} synchrone, donc exécutées dans le threadpool (SSH bloquant)\\
Backend & \fn{kpi\_service.get\_overview} : \fn{\_storage\_health()} sert le cache (\(\le 60\)\,s), sinon rafraîchit\\
Tool & \fn{get\_storage\_health()} --- \textbf{appel direct, pas via MCP} --- lance \fn{get\_capacity} et \fn{list\_backups} en parallèle dans un \cls{ThreadPoolExecutor}\\
Infrastructure & Deux sessions SSH : \texttt{df -T -B1 -P} et \texttt{systemctl list-timers} + \texttt{systemctl show}\\
Résultat & Capacité, jobs, alertes déterministes (seuils 90/95\,\% et 24\,h), \texttt{overall\_status}\\
Backend (suite) & Agrégats SQL sur \tbl{incident\_tickets}, calcul du \texttt{health\_score} et du \texttt{headline}\\
Frontend & \cls{KpiCard}, \cls{CapacityOverview}, \cls{PerformanceChart}, \cls{RecentActivity}, \cls{OasisInsights}, \cls{AttentionBanner}\\
\bottomrule
\end{longtable}

\begin{keybox}
\textbf{Pourquoi ce workflow ne passe pas par MCP ni par le LLM.} Le dashboard
doit s'afficher vite et rendre toujours la même valeur pour le même état.
Passer par le sous-processus MCP ajouterait une latence de démarrage, et passer
par le LLM rendrait l'affichage non déterministe. Le fichier le déclare :
« Aucun appel LLM ici. »
\end{keybox}

% =====================================================================
\section{Workflow 3 --- Chat avec l'assistant}

\begin{figure}[h!]
\centering
\begin{tikzpicture}[node distance=6mm]
\node[box,minimum width=100mm] (u) {\textbf{Utilisateur} --- « Quels volumes approchent de leur limite ? »};
\node[bluebox,below=of u,minimum width=100mm] (f) {\textbf{ChatPage} \quad api.post('/storage/chat', \{message, conversation\_id\})};
\node[bluebox,below=of f,minimum width=100mm] (a) {\textbf{API} \quad require\_role \arrowto{} history\_for\_llm(10 msg) \arrowto{} contexte incident éventuel};
\node[vibox,below=of a,minimum width=100mm] (ag) {\textbf{StoragePersona.handle\_message} \quad RAG Chroma + prompt système + historique};
\node[vibox,below=of ag,minimum width=100mm] (g1) {\textbf{Groq} \quad chat\_with\_tools(messages, tools) \arrowto{} tool\_calls = [get\_capacity]};
\node[ambox,below=of g1,minimum width=100mm] (m) {\textbf{MCP} \quad client stdio \arrowto{} MCPServer \arrowto{} get\_capacity()};
\node[redbox,below=of m,minimum width=100mm] (v) {\textbf{VM} \quad ssh: df -T -B1 -P \arrowto{} \{"volumes": [\ldots]\}};
\node[greenbox,below=of v,minimum width=100mm] (au) {\textbf{Audit} \quad INSERT mcp\_calls (tool\_name, arguments, result, correlation\_id)};
\node[vibox,below=of au,minimum width=100mm] (g2) {\textbf{Groq} \quad second tour : rédige la réponse à partir du résultat};
\node[bluebox,below=of g2,minimum width=100mm] (s) {\textbf{save\_turn} \quad INSERT 2 chat\_messages + updated\_at};
\node[box,below=of s,minimum width=100mm,draw=oasisdark] (r) {\textbf{ChatTurnOut} \arrowto{} affichage React};
\foreach \x/\y in {u/f,f/a,a/ag,ag/g1,g1/m,m/v,v/au,au/g2,g2/s,s/r} \draw[flowb] (\x)--(\y);
\end{tikzpicture}
\caption{Un tour de chat avec un outil de lecture.}
\end{figure}

\clearpage

\subsection*{Variante : le modèle veut relancer une sauvegarde}
À l'étape MCP, le nom appartient à \texttt{REMEDIATION\_TOOLS}. L'agent
n'exécute rien : il crée une \tbl{action\_requests} \texttt{pending} et renvoie
au modèle un résultat qui \emph{lui dit} de présenter la chose comme une
recommandation. Le tour se termine par un message accompagné d'un composant
\cls{ActionProposal}.

\subsection*{Effet de bord : création automatique de ticket}
Si \fn{evaluate\_tool\_result} détecte une anomalie (volume \(\ge\) 85\,\%, job
en échec), un \tbl{incident\_tickets} est créé à la fin du tour, attribué à
l'utilisateur courant, avec le \texttt{correlation\_id} du tour --- ce qui relie
le ticket aux \tbl{mcp\_calls} qui l'ont motivé.

% =====================================================================
\section{Workflow 4 --- Détection automatique d'un incident}

\textbf{Aucun utilisateur n'est impliqué.} C'est une tâche de fond démarrée par
le \fn{lifespan} de \file{main.py}.

\begin{longtable}{@{}p{2.9cm}p{11.9cm}@{}}
\toprule
\textbf{Étape} & \textbf{Ce qui se passe}\\
\midrule
\endhead
Déclencheur & \fn{incident\_watcher.run\_forever}, boucle toutes les 300\,s\\
Isolation & \fn{asyncio.to\_thread(scan\_once)} --- la lecture SSH est bloquante et ne doit pas figer la boucle d'événements\\
Tool & \fn{list\_backups()} en appel direct\\
Infrastructure & \texttt{systemctl list-timers -{}-all}, puis \texttt{systemctl show} sur toutes les units \textbf{en une seule commande} (séparateur \texttt{\#\#\#OASIS-UNIT=})\\
Règle & \fn{check\_backup\_status(job)} : seulement si \texttt{status == "failed"}\\
Vérification d'abord & \fn{\_verify\_resolved} referme les tickets dont une action \texttt{completed} existe et dont le job est redevenu sain\\
Anti-doublon & \fn{\_already\_handled} : on ne rouvre qu'à partir d'une exécution postérieure au ticket précédent\\
Backend & \fn{create\_ticket\_from\_diagnosis} avec \texttt{created\_by = compte de service}\\
Audit & Première note : « Detected automatically by monitoring: job \ldots{} came back failed (systemd result: \ldots) »\\
Frontend & Le ticket apparaît au prochain chargement de \route{/storage/incidents}\\
\bottomrule
\end{longtable}

\begin{infobox}
\textbf{Pourquoi cette tâche existe.} Le docstring le justifie : « les tickets ne
sont créés que par le persona Storage, pendant une conversation. Un job qui
échoue la nuit ne produit donc un ticket que si quelqu'un pense à interroger
l'agent le lendemain. Les alertes du dashboard, elles, sont recalculées à chaque
appel et ne laissent aucune trace. »
\end{infobox}

% =====================================================================
\section{Workflow 5 --- Investigation d'un incident}

\begin{longtable}{@{}p{2.9cm}p{11.9cm}@{}}
\toprule
\textbf{Étape} & \textbf{Ce qui se passe}\\
\midrule
\endhead
Utilisateur & Ouvre un ticket, clique « Investigate with OASIS »\\
Frontend & \route{POST /incidents/\{id\}/investigate}\\
API & \fn{require\_role("storage\_engineer","administrator")}. Route \texttt{async}, volontairement bloquante : « l'ingénieur vient de cliquer et attend le diagnostic, un accusé de réception vide ne lui apprendrait rien »\\
Backend & \fn{incident\_investigation.start} : fil existant \arrowto{} retour immédiat, \texttt{created=False}\\
Agent & \fn{handle\_message} avec \cls{IncidentContext} : le prompt contient le numéro du ticket, la cible technique, le relevé MCP et trois consignes\\
MCP & Le modèle appelle typiquement \fn{get\_backup\_logs} puis \fn{get\_backup\_job\_status}\\
Infrastructure & \texttt{journalctl -u <service> -n 20 -o short-iso}, \texttt{systemctl show}\\
Proposition & Le modèle appelle un outil de remédiation \arrowto{} \tbl{action\_requests} \texttt{pending}\\
Filet & Si aucune proposition n'a été enregistrée : \fn{\_recover\_written\_call} puis \fn{\_fallback\_action}\\
Écriture & Conversation créée, \texttt{incident\_ticket\_id} posé, \fn{save\_turn}, \fn{attach\_actions}, \fn{mark\_in\_progress}\\
Synthèse & Second appel Groq \arrowto{} JSON \texttt{\{root\_cause, impact\_summary\}} recopié sur le ticket\\
Audit & Note « Investigation opened by \ldots{} proposed N action(s) for confirmation »\\
Frontend & \cls{InvestigationOut} \arrowto{} ouverture du fil, ticket rechargé avec sa fiche complète\\
\bottomrule
\end{longtable}

\begin{keybox}
\textbf{Trois filets successifs pour ne pas perdre la proposition.}
\emph{(1)} l'appel d'outil normal ;
\emph{(2)} \fn{\_recover\_written\_call}, qui récupère un appel écrit en JSON
dans la prose du modèle ;
\emph{(3)} \fn{\_fallback\_action}, qui reprend l'\texttt{action\_type} calculé
par la règle de détection au moment de l'ouverture du ticket.
Le cycle ne doit pas « s'arrêter à l'avant-dernière marche ».
\end{keybox}

% =====================================================================
\section{Workflow 6 --- Confirmation humaine et clôture}

C'est le workflow central du projet. Il répond exactement à la chaîne
\emph{demande \arrowto{} analyse \arrowto{} proposition \arrowto{} confirmation
\arrowto{} exécution \arrowto{} résultat}.

\begin{figure}[h!]
\centering
\begin{tikzpicture}[node distance=5.5mm]
\node[box,minimum width=108mm] (n1) {\textbf{1. Analyse} --- l'agent lit la VM et identifie la cause racine};
\node[ambox,below=of n1,minimum width=108mm] (n2) {\textbf{2. Proposition} --- \_propose\_action() \arrowto{} action\_requests(status=pending), \textbf{rien n'est exécuté}};
\node[bluebox,below=of n2,minimum width=108mm] (n3) {\textbf{3. Affichage} --- ActionProposal : libellé lisible + action\_type(target) + avertissement si destructif};
\node[greenbox,below=of n3,minimum width=108mm] (n4) {\textbf{4. Décision humaine} --- clic « Confirm » ou « Reject »};
\node[ambox,below=of n4,minimum width=108mm] (n5) {\textbf{5. Démarrage MCP} --- StorageMCPClient().start() ; échec \arrowto{} HTTP 503, action toujours PENDING};
\node[redbox,below=of n5,minimum width=108mm] (n6) {\textbf{6. Exécution} --- call\_tool(action\_type, confirm=True, **parameters) \arrowto{} SSH \arrowto{} VM};
\node[greenbox,below=of n6,minimum width=108mm] (n7) {\textbf{7. Qualification} --- result contient "error" ? FAILED : COMPLETED ; audit dans mcp\_calls};
\node[bluebox,below=of n7,minimum width=108mm] (n8) {\textbf{8. Vérification} --- incident\_verification.verify() relit l'indicateur sur la VM};
\node[box,below=of n8,minimum width=108mm,draw=oasisdark] (n9) {\textbf{9. Clôture} --- si verdict.resolved : resolve\_ticket + note ; sinon le ticket reste ouvert avec la raison};
\foreach \x/\y in {n1/n2,n2/n3,n3/n4,n4/n5,n5/n6,n6/n7,n7/n8,n8/n9} \draw[flowb] (\x)--(\y);
\end{tikzpicture}
\caption{Les neuf temps de la confirmation humaine, tels qu'ils sont codés.}
\end{figure}

\clearpage

\subsection*{Ce que voit l'ingénieur avant de cliquer}
\cls{ActionProposal} traduit l'appel d'outil en phrase :

\begin{lstlisting}[caption={frontend/src/components/ActionProposal.tsx}]
const ACTION_LABELS = {
  run_backup:               (t) => `Re-run backup job ${t} now`,
  create_snapshot:          (t) => `Take a snapshot of volume ${t}`,
  restore_from_backup:      (t) => `Restore ${t} from a backup`,
  repair_backup_repository: (t) => `Repair the backup repository configured for ${t}`,
};
const DESTRUCTIVE_ACTIONS = new Set(['restore_from_backup']);
\end{lstlisting}

\noindent
Sous la phrase, l'appel technique brut est affiché en police à chasse fixe :
\texttt{action\_type(target)}. Pour une action destructive, un avertissement
rouge s'ajoute : « This action overwrites existing data. Check the target before
confirming. »

\subsection*{La branche « Reject »}
\route{POST /storage/actions/\{id\}/reject} passe l'action en \texttt{REJECTED}
et, si elle porte un ticket, écrit une note nommée avec le motif éventuel :
« Un rejet est une décision, pas un non-événement : il doit rester lisible dans
l'historique du ticket. »

% =====================================================================
\section{Workflow 7 --- Data Hub et RAG}

\begin{longtable}{@{}p{2.9cm}p{11.9cm}@{}}
\toprule
\textbf{Étape} & \textbf{Ce qui se passe}\\
\midrule
\endhead
Utilisateur & Dépose un fichier sur \route{/storage/data-hub}\\
Frontend & \fn{api.upload('/storage/documents', file)} --- \texttt{FormData}\\
API & \route{POST /storage/documents}, \fn{require\_role}, \texttt{await file.read()}\\
Backend & \fn{knowledge\_service.create\_document} : taille \(\le\) 10\,Mo, extension autorisée, écriture sous \texttt{<uuid><ext>}\\
Extraction & Texte brut, ou \texttt{pypdf} pour un PDF ; refus si \(<\) 20 mots utiles\\
Découpage & \fn{chunk\_text} : 500 mots, recouvrement 50\\
Vectorisation & \fn{embed\_passages} avec le préfixe \texttt{"passage: "}\\
Chroma & \texttt{collection.upsert} dans \texttt{storage\_kb}, identifiants dérivés de l'UUID\\
Résultat & \texttt{status=indexed} et \texttt{chunk\_count}, ou \texttt{status=failed} avec le message d'erreur conservé\\
Frontend & La ligne apparaît avec son statut ; un document en échec explique pourquoi\\
Plus tard & Chaque question du chat déclenche \fn{embed\_query} (préfixe \texttt{"query: "}) et les 3 passages les plus proches entrent dans le prompt\\
\bottomrule
\end{longtable}

\begin{keybox}
\textbf{Pourquoi deux fonctions d'embedding et non une seule.} Les modèles E5
sont \emph{asymétriques} : entraînés avec \texttt{"query: "} devant la question
et \texttt{"passage: "} devant le document. Le fichier le dit : « une fonction
unique \fn{embed()} serait forcément fausse d'un côté ou de l'autre ».
\end{keybox}

% =====================================================================
\section{Workflow 8 --- Demandes de service}

\begin{longtable}{@{}p{2.9cm}p{11.9cm}@{}}
\toprule
\textbf{Étape} & \textbf{Ce qui se passe}\\
\midrule
\endhead
Utilisateur & Ingénieur Storage : « J'ai besoin de 500\,Go supplémentaires »\\
Frontend & \route{POST /requests} avec titre, description, type, priorité\\
API & \fn{require\_role("storage\_engineer","administrator")} --- le Linux Engineer n'a pas accès à ce canal\\
Backend & \fn{create\_request} : la \textbf{direction} est déduite du rôle, jamais du corps de la requête\\
Destinataire & S'il n'y a qu'un administrateur actif, il est choisi automatiquement ; s'il y en a plusieurs et qu'aucun n'est précisé, la demande est refusée\\
Base & \texttt{INSERT service\_requests} (\texttt{status=pending}) + \texttt{INSERT request\_events} (\texttt{to\_status=pending}, note « Demande créée »)\\
Administrateur & Voit la demande dans sa boîte « received » (\route{GET /requests?box=received})\\
Transition & \route{PATCH /requests/\{id\}/status} : contrôle \texttt{assigned\_to == user.id}, puis table \fn{ALLOWED\_TRANSITIONS}\\
Décision & \texttt{approved} ou \texttt{rejected} \arrowto{} \texttt{decided\_at} figé, \texttt{administrator\_response} enregistrée\\
Audit & Une ligne \tbl{request\_events} par transition, jamais renvoyée au client\\
Frontend & Le statut et la date de décision s'affichent dans la liste et le détail\\
\bottomrule
\end{longtable}

\noindent
\textbf{Aucun LLM, aucun MCP, aucune infrastructure} n'intervient : c'est un
canal humain.

% =====================================================================
\section{Workflow 9 --- Recherche globale}

Utilisateur \arrowto{} \cls{GlobalSearch} (Topbar) \arrowto{}
\route{GET /storage/search?q=} \arrowto{} \fn{search\_service.search}
\arrowto{} quatre requêtes \texttt{ILIKE} échappées (incidents, demandes,
conversations, documents) \arrowto{} liste de \cls{SearchHitOut}
(\texttt{type}, \texttt{id}, \texttt{title}, \texttt{subtitle},
\texttt{timestamp}) \arrowto{} le frontend traduit le \texttt{type} en route :

\begin{lstlisting}[caption={frontend/src/components/layout/GlobalSearch.tsx}]
const DESTINATION = {
  incident:     (id) => `/storage/incidents?incident=${id}`,
  request:      (id) => `/storage/requests?request=${id}`,
  conversation: (id) => `/storage/chat?conversation=${id}`,
  document:     (id) => `/storage/data-hub?document=${id}`,
};
\end{lstlisting}

% =====================================================================
\section{Workflow 10 --- Rapports PDF}

Utilisateur \arrowto{} \cls{ReportsPage} \arrowto{} \fn{api.download} \arrowto{}
\route{GET /storage/reports/global} ou
\route{GET /storage/reports/incidents/\{id\}} \arrowto{}
\fn{report\_service.build\_*\_report} (lecture SQL + lecture de capacité par SSH)
\arrowto{} \texttt{bytes} PDF \arrowto{} \texttt{Response} avec
\texttt{Content-Disposition} et \texttt{Access-Control-Expose-Headers}
\arrowto{} le frontend lit le blob, extrait le nom de fichier et déclenche le
téléchargement.

% =====================================================================
\section{Workflow 11 --- Accès cross-domain}

\begin{enumerate}
  \item \route{POST /cross-domain/requests} --- n'importe quel utilisateur
  connecté crée une demande (\texttt{target\_workspace}, \texttt{reason}).
  \item \route{GET /cross-domain/requests/pending} --- réservé à
  \texttt{administrator}.
  \item \route{POST /cross-domain/requests/\{id\}/decision} --- l'admin approuve
  ou refuse. En cas d'approbation, \fn{create\_cross\_domain\_token} génère un
  JWT de 30 minutes avec \texttt{scope="cross\_domain\_read\_only"} et le
  workspace, stocké dans \texttt{granted\_token}.
\end{enumerate}

\begin{warnbox}
\textbf{Le workflow s'arrête là.} Aucun code ne vérifie ensuite ce scope, et
aucune interface frontend ne l'expose. C'est une brique préparée et testable,
pas une fonctionnalité utilisable de bout en bout aujourd'hui.
\end{warnbox}

% =====================================================================
\section{Workflow 12 --- Collecte de métriques}

\begin{longtable}{@{}p{2.9cm}p{11.9cm}@{}}
\toprule
\textbf{Étape} & \textbf{Ce qui se passe}\\
\midrule
\endhead
Déclencheur & \fn{metrics\_collector.run\_forever}, toutes les 300\,s\\
Isolation & \fn{asyncio.to\_thread(\_collect)}\\
Tool & \fn{sample\_disk\_performance()} --- \textbf{appel direct} ; cet outil n'est \emph{pas} dans le \fn{TOOL\_REGISTRY} et n'est donc pas exposé au LLM\\
Infrastructure & Une session SSH : deux lectures de \texttt{/proc/diskstats} espacées de 2\,s\\
Calcul & \fn{\_rates} : IOPS, débit (secteurs \(\times\) 512), latence moyenne ; \texttt{None} si un compteur a reculé\\
Base & Une ligne \tbl{disk\_metrics} par device\\
Purge & Tous les 288 cycles (\(\approx\) une fois par jour), suppression au-delà de 31 jours\\
Lecture & \route{GET /storage/dashboard/performance?range=} \arrowto{} \fn{read\_series} (\texttt{date\_bin}, agrégation en deux passes)\\
Frontend & \cls{PerformanceChart} (recharts) ; si \texttt{collecting} est vrai, un message remplace la courbe\\
\bottomrule
\end{longtable}

% =====================================================================
\section{Synthèse : quel workflow mobilise quoi ?}

\begin{longtable}{@{}p{3.6cm}ccccc@{}}
\toprule
\textbf{Workflow} & \textbf{LLM} & \textbf{RAG} & \textbf{MCP} & \textbf{SSH} & \textbf{Confirmation}\\
\midrule
\endhead
Authentification & --- & --- & --- & --- & ---\\
Dashboard & --- & --- & --- & \(\bullet\) & ---\\
Chat & \(\bullet\) & \(\bullet\) & \(\bullet\) & \(\bullet\) & si remédiation\\
Détection auto & --- & --- & --- & \(\bullet\) & ---\\
Investigation & \(\bullet\) & \(\bullet\) & \(\bullet\) & \(\bullet\) & \(\bullet\)\\
Confirmation & --- & --- & \(\bullet\) & \(\bullet\) & \(\bullet\)\\
Data Hub & --- & \(\bullet\) & --- & --- & ---\\
Demandes & --- & --- & --- & --- & ---\\
Recherche & --- & --- & --- & --- & ---\\
Rapports & --- & --- & --- & \(\bullet\) & ---\\
Cross-domain & --- & --- & --- & --- & ---\\
Métriques & --- & --- & --- & \(\bullet\) & ---\\
\bottomrule
\end{longtable}
